% ===== 09_conclusion =====
\section{结论}
\label{sec:conclusion}
接触控制与追踪隔离的组合不改变阈值策略的进入与退出状态，只改变其间的路径；控制时长把这条路径上的速度、感染与未感染者的隔离联系在一起。控制越快，感染越少，但被隔离的未感染者越多，因此不同目标对应不同的最优控制。最快、感染最少的控制以追踪隔离为主，隔离负担最小的控制以减少接触为主，成本最优控制则在控制后期让两种措施分担干预强度，居于两者之间。在群体免疫阈值退出避免了疫情反弹，却给累计新增感染设定了与控制组合无关的下界。因此，设计沿阈值运行的干预时，既要选择两种干预的组合，也要把退出要求纳入评价。
